\documentclass[a4paper,10pt,twocolumn]{article}
\usepackage[utf8]{inputenc}
\usepackage[T1]{fontenc}
\usepackage[indonesian]{babel}
\usepackage{lmodern}
\usepackage{geometry}
\usepackage{amsmath,amssymb}
\usepackage{booktabs}
\usepackage{array}
\usepackage{enumitem}
\usepackage{xcolor}
\usepackage{listings}
\usepackage{fancyhdr}
\usepackage{titlesec}
\usepackage{microtype}

\geometry{top=1.35cm,bottom=1.45cm,left=1.25cm,right=1.25cm,columnsep=0.65cm}
\setlength{\parindent}{0pt}
\setlength{\parskip}{3pt}
\setlist[itemize]{leftmargin=*, itemsep=1pt, topsep=2pt}
\setlist[enumerate]{leftmargin=*, itemsep=2pt, topsep=2pt}
\titleformat{\section}{\large\bfseries\color{black}}{}{0pt}{}
\titleformat{\subsection}{\normalsize\bfseries\color{black}}{}{0pt}{}
\pagestyle{fancy}
\fancyhf{}
\lhead{Latihan UAS Machine Learning - Exercise 2}
\rhead{Ringkasan Jawaban}
\cfoot{\thepage}

\lstset{
    basicstyle=\ttfamily\scriptsize,
    breaklines=true,
    frame=single,
    columns=fullflexible,
    keepspaces=true,
    showstringspaces=false,
    xleftmargin=0pt,
    xrightmargin=0pt
}

\newcommand{\soal}[1]{\textbf{Soal.} #1}
\newcommand{\jawab}{\textbf{Jawaban singkat.} }
\newcommand{\catatan}[1]{\textit{Catatan: #1}}

\begin{document}

\begin{center}
    {\Large \textbf{Pembahasan Latihan Soal UAS Machine Learning}}\\[-1pt]
    {\large \textbf{Exercise Machine Learning - 2}}\\
    Lecturer: Afiahayati, S.Kom., M.Cs., Ph.D.\\
    \vspace{2pt}
    \small Format bahan belajar: soal + jawaban ringkas, 2 kolom per halaman.
\end{center}

\section{K-Nearest Neighbor (3-NN)}

\soal{Lihat Tabel 1. Tentukan kelas $(Y)$ untuk data uji menggunakan metode 3-nearest neighbor dengan jarak Euclidean.}

\begin{center}
\scriptsize
\begin{tabular}{c|cccccc|cc}
\toprule
 & \multicolumn{6}{c|}{Training Data} & \multicolumn{2}{c}{Testing Data} \\
\midrule
$X_1$ & 2 & 1 & 2 & 3 & 1 & 3 & 2 & 1 \\
$X_2$ & 2 & 2 & 3 & 3 & 0 & 2 & 1 & 1 \\
$Y$   & -1 & 1 & -1 & 1 & 1 & -1 & ? & ? \\
\bottomrule
\end{tabular}
\end{center}

\jawab
Rumus jarak Euclidean:
\[
d(p,q)=\sqrt{(x_{1p}-x_{1q})^2+(x_{2p}-x_{2q})^2}.
\]

Misalkan data latih:
\[
\begin{array}{c|c|c}
\text{ID} & (X_1,X_2) & Y \\
\hline
A&(2,2)&-1\\
B&(1,2)&1\\
C&(2,3)&-1\\
D&(3,3)&1\\
E&(1,0)&1\\
F&(3,2)&-1
\end{array}
\]

\textbf{Data uji 1:} $T_1=(2,1)$.

Hitung jarak $T_1$ ke semua data latih:
\[
\begin{aligned}
d(T_1,A)&=\sqrt{(2-2)^2+(1-2)^2}=1,\\
d(T_1,B)&=\sqrt{(2-1)^2+(1-2)^2}=\sqrt2,\\
d(T_1,C)&=\sqrt{(2-2)^2+(1-3)^2}=2,\\
d(T_1,D)&=\sqrt{(2-3)^2+(1-3)^2}=\sqrt5,\\
d(T_1,E)&=\sqrt{(2-1)^2+(1-0)^2}=\sqrt2,\\
d(T_1,F)&=\sqrt{(2-3)^2+(1-2)^2}=\sqrt2.
\end{aligned}
\]

\begin{center}
\scriptsize
\begin{tabular}{c|cccccc}
\toprule
ID & A & B & C & D & E & F \\
\midrule
$d(T_1,\cdot)$ & $1$ & $\sqrt2$ & $2$ & $\sqrt5$ & $\sqrt2$ & $\sqrt2$ \\
$Y$ & -1 & 1 & -1 & 1 & 1 & -1 \\
\bottomrule
\end{tabular}
\end{center}

Urutan jarak terkecil: $A(1)$, lalu seri $B,E,F(\sqrt2)$, kemudian $C(2)$ dan $D(\sqrt5)$. Dengan aturan tie-break urutan data, tiga tetangga terdekat adalah $A(-1),B(1),E(1)$. Voting: dua kelas $1$ dan satu kelas $-1$, sehingga mayoritas = $1$.
\[
\boxed{Y(T_1)=1}.
\]
\catatan{Pada $T_1$ ada jarak seri di $\sqrt2$. Jika dosen menentukan aturan tie-break lain, tetangga ketiga bisa berubah. Di sini dipakai aturan urutan data karena soal tidak memberi aturan khusus.}

\textbf{Data uji 2:} $T_2=(1,1)$.

Hitung jarak $T_2$ ke semua data latih:
\[
\begin{aligned}
d(T_2,A)&=\sqrt{(1-2)^2+(1-2)^2}=\sqrt2,\\
d(T_2,B)&=\sqrt{(1-1)^2+(1-2)^2}=1,\\
d(T_2,C)&=\sqrt{(1-2)^2+(1-3)^2}=\sqrt5,\\
d(T_2,D)&=\sqrt{(1-3)^2+(1-3)^2}=\sqrt8,\\
d(T_2,E)&=\sqrt{(1-1)^2+(1-0)^2}=1,\\
d(T_2,F)&=\sqrt{(1-3)^2+(1-2)^2}=\sqrt5.
\end{aligned}
\]

\begin{center}
\scriptsize
\begin{tabular}{c|cccccc}
\toprule
ID & A & B & C & D & E & F \\
\midrule
$d(T_2,\cdot)$ & $\sqrt2$ & $1$ & $\sqrt5$ & $\sqrt8$ & $1$ & $\sqrt5$ \\
$Y$ & -1 & 1 & -1 & 1 & 1 & -1 \\
\bottomrule
\end{tabular}
\end{center}

Urutan jarak terkecil: $B(1)$, $E(1)$, $A(\sqrt2)$, lalu $C,F(\sqrt5)$ dan $D(\sqrt8)$. Tiga tetangga terdekat: $B(1),E(1),A(-1)$. Voting: dua kelas $1$ dan satu kelas $-1$, sehingga mayoritas = $1$.
\[
\boxed{Y(T_2)=1}.
\]

\textbf{Hasil akhir testing data:}
\[
\boxed{(2,1)\rightarrow 1, \qquad (1,1)\rightarrow 1}
\]

\section{Support Vector Machine}

\soal{Jelaskan untuk data seperti apa Support Vector Machine cocok digunakan.}

\jawab
SVM cocok untuk:
\begin{itemize}
    \item Data klasifikasi numerik yang dapat dipisahkan oleh hyperplane, baik linear maupun non-linear.
    \item Data berdimensi tinggi, misalnya teks, fitur vektor, atau citra yang sudah diekstraksi fiturnya.
    \item Dataset kecil sampai menengah, terutama saat ingin margin pemisah yang besar.
    \item Data non-linear jika memakai kernel, misalnya RBF atau polynomial.
\end{itemize}
Sebelum memakai SVM, fitur sebaiknya dinormalisasi atau distandarisasi agar skala fitur tidak mendominasi jarak/margin.

\section{Naive Bayes: CategoricalNB}

\soal{Perhatikan kode pada Gambar 1. Jelaskan cara kerja metode machine learning pada kode tersebut.}

\textbf{Deskripsi Gambar 1.} Gambar berisi kode Python yang mengimpor modul \texttt{naive\_bayes}, membuat model \texttt{CategoricalNB()}, melatih model dengan \texttt{fit(X\_train, y\_train)}, lalu memprediksi data uji dengan \texttt{predict(X\_test)}.

\begin{lstlisting}[language=Python]
from sklearn import naive_bayes
nb = naive_bayes.CategoricalNB()
nb.fit(X_train, y_train)
y_pred = nb.predict(X_test)
\end{lstlisting}

\jawab
Kode tersebut menggunakan \textbf{Naive Bayes untuk fitur kategorikal}. Cara kerjanya:
\begin{enumerate}
    \item \texttt{CategoricalNB()} membuat model Naive Bayes untuk data kategori/diskrit.
    \item \texttt{fit(X\_train, y\_train)} menghitung prior kelas $P(Y)$ dan peluang fitur kategori terhadap kelas $P(X_i \mid Y)$.
    \item Model mengasumsikan setiap fitur saling independen bersyarat terhadap kelas.
    \item \texttt{predict(X\_test)} memilih kelas dengan posterior terbesar:
\end{enumerate}
\[
\hat{y}=\arg\max_c P(c)\prod_{i=1}^{d}P(x_i\mid c).
\]
Hasil akhirnya disimpan pada \texttt{y\_pred}.

\section{Ensemble Learning}

\subsection{Bagging, Boosting, dan Stacking}

\soal{Jelaskan algoritma bagging, boosting, dan stacking serta perbedaannya.}

\jawab
\textbf{Bagging} atau bootstrap aggregating:
\begin{itemize}
    \item Membuat beberapa dataset baru dari data latih dengan sampling bootstrap.
    \item Melatih beberapa model secara paralel.
    \item Prediksi akhir diperoleh dari voting mayoritas atau rata-rata.
    \item Tujuan utama: menurunkan variance dan mengurangi overfitting.
\end{itemize}
Contoh: Random Forest.

\textbf{Boosting}:
\begin{itemize}
    \item Melatih model lemah secara berurutan.
    \item Model berikutnya lebih fokus pada data yang salah diprediksi oleh model sebelumnya.
    \item Prediksi akhir berupa voting berbobot.
    \item Tujuan utama: menurunkan bias dan meningkatkan akurasi.
\end{itemize}
Contoh: AdaBoost, Gradient Boosting.

\textbf{Stacking}:
\begin{itemize}
    \item Melatih beberapa model dasar yang bisa berbeda jenis.
    \item Output model dasar dijadikan input untuk meta-model.
    \item Meta-model belajar menggabungkan prediksi model dasar.
\end{itemize}

\textbf{Perbedaan inti:}
\begin{center}
\tiny
\begin{tabular}{p{1.5cm}|p{1.55cm}|p{1.55cm}|p{1.55cm}}
\toprule
Aspek & Bagging & Boosting & Stacking \\
\midrule
Pola latih & Paralel & Berurutan & Dua level \\
Data & Bootstrap & Berbobot, fokus error & Prediksi base model \\
Gabungan & Voting/rata-rata & Voting berbobot & Meta-model \\
Fokus & Kurangi variance & Kurangi bias & Kombinasi model \\
\bottomrule
\end{tabular}
\end{center}

\subsection{Parameter AdaBoost}

\soal{Perhatikan Gambar 2. Jelaskan arti \texttt{base\_classifier}, \texttt{n\_estimators}, dan \texttt{learning\_rate}.}

\textbf{Deskripsi Gambar 2.} Gambar berisi potongan kode pembuatan \texttt{AdaBoostClassifier} dengan parameter \texttt{base\_classifier}, \texttt{n\_estimators=50}, \texttt{learning\_rate=1.0}, dan \texttt{random\_state=42}; lalu model dilatih menggunakan \texttt{fit(X\_train, y\_train)}.

\begin{lstlisting}[language=Python]
adaboost_classifier = AdaBoostClassifier(
    base_classifier,
    n_estimators=50,
    learning_rate=1.0,
    random_state=42
)
adaboost_classifier.fit(X_train, y_train)
\end{lstlisting}

\jawab
\begin{itemize}
    \item \texttt{base\_classifier}: model dasar/weak learner yang dipakai AdaBoost, misalnya decision stump.
    \item \texttt{n\_estimators=50}: jumlah weak learner yang akan dibuat dalam proses boosting.
    \item \texttt{learning\_rate=1.0}: mengatur besar kontribusi setiap weak learner terhadap prediksi akhir. Nilai kecil membuat pembelajaran lebih hati-hati.
    \item \texttt{random\_state=42}: membuat hasil random dapat direproduksi.
\end{itemize}

\subsection{Implementasi Bagging Manual}

\soal{Perhatikan Tabel 2. Implementasikan algoritma Bagging untuk data dengan hanya 3 bootstrapping/single model. Tulis prediksi akhir dan hitung akurasinya.}

\begin{center}
\scriptsize
\begin{tabular}{c|cccccc}
\toprule
$x$ & 2.4 & 3.1 & 3.7 & 4.4 & 5.2 & 5.8 \\
$y$ & 1 & 1 & 1 & -1 & -1 & -1 \\
\bottomrule
\end{tabular}
\end{center}

\jawab
\catatan{Soal tidak menentukan sampel bootstrap acak dan jenis base learner. Untuk latihan manual, digunakan asumsi base learner berupa decision stump/threshold pada fitur $x$. Karena data terpisah sempurna antara $x=3.7$ dan $x=4.4$, batas yang wajar adalah $x=4.05$.}

Aturan model tunggal:
\[
h(x)=
\begin{cases}
1, & x < 4.05,\\
-1, & x \ge 4.05.
\end{cases}
\]

Contoh 3 bootstrap dan model yang terbentuk:
\begin{center}
\scriptsize
\begin{tabular}{c|p{2.3cm}|c}
\toprule
Model & Indeks bootstrap & Aturan \\
\midrule
$h_1$ & 1,2,3,4,5,5 & $x<4.05\Rightarrow 1$ \\
$h_2$ & 1,1,2,3,4,6 & $x<4.05\Rightarrow 1$ \\
$h_3$ & 2,3,3,4,5,6 & $x<4.05\Rightarrow 1$ \\
\bottomrule
\end{tabular}
\end{center}

Prediksi masing-masing model pada data sumber:
\begin{center}
\scriptsize
\begin{tabular}{c|cccccc}
\toprule
$x$ & 2.4 & 3.1 & 3.7 & 4.4 & 5.2 & 5.8 \\
\midrule
$y$ asli & 1 & 1 & 1 & -1 & -1 & -1 \\
$h_1$ & 1 & 1 & 1 & -1 & -1 & -1 \\
$h_2$ & 1 & 1 & 1 & -1 & -1 & -1 \\
$h_3$ & 1 & 1 & 1 & -1 & -1 & -1 \\
Voting & 1 & 1 & 1 & -1 & -1 & -1 \\
\bottomrule
\end{tabular}
\end{center}

Prediksi akhir Bagging:
\[
\boxed{\hat{y}=[1,1,1,-1,-1,-1]}.
\]

Akurasi:
\[
\text{Accuracy}=\frac{\text{jumlah benar}}{\text{jumlah data}}=\frac{6}{6}=1=\boxed{100\%}.
\]

\section*{Ringkasan Jawaban Cepat}
\begin{itemize}
    \item 3-NN testing: $(2,1)\rightarrow 1$ dan $(1,1)\rightarrow 1$ dengan tie-break urutan data.
    \item SVM cocok untuk data numerik, berdimensi tinggi, dan memiliki margin pemisah yang jelas; bisa non-linear dengan kernel.
    \item \texttt{CategoricalNB}: menghitung prior dan conditional probability kategori, lalu memilih kelas dengan posterior terbesar.
    \item Bagging: paralel + bootstrap + voting. Boosting: sequential + fokus error + voting berbobot. Stacking: base model + meta-model.
    \item AdaBoost: \texttt{base\_classifier} = weak learner, \texttt{n\_estimators} = jumlah learner, \texttt{learning\_rate} = bobot kontribusi learner.
    \item Bagging manual pada Tabel 2: prediksi $[1,1,1,-1,-1,-1]$, akurasi $100\%$ berdasarkan asumsi stump $x<4.05$.
\end{itemize}

\end{document}
